\subsection{Relative geometry and regular finite observation}
The forward mechanism has a precise local analytic comparator.
Section~\ref{sec:v14-normal-form} derives the relative endpoint flux
from a hyperbolic normal form after the necessary projection to
physical endpoints. Its individually normalized factors agree, in
the analytic overlap, with the half-line determinant amplitudes.
Thus the existence of an analytic relative product is not used as a
stand-alone novelty claim. The full smooth statement joins the
explicit determinant construction, mixed derivative control, physical
channel and residual-time integration, and the complete-profile
inverse and observation theorem.

For finite-dimensional observation the relevant hypothesis is regular
rank, not dimension by itself. Theorem~\ref{thm:v14-rank} characterizes
local bi-Lipschitz observation by positive windows on a supplied
family. Theorem~\ref{thm:v13-two-flight-observation} verifies this for
the constructed fixed-area contact families. In its raw-probability
experiment the supplied family includes $A$, whereas the preceding
geometric inverse may instead take normalized exact germs as data.
Section~\ref{sec:v14-normalization} makes this distinction explicit.

There is also a direct finite-window experiment with unknown area.
Keeping the existing area compensator free gives a physical family
with coordinates $(A,q_2,\ldots,q_M)$ at fixed labelled
$g,\kappa_0,\kappa_1$. Theorem~\ref{thm:v14-area-windows} recovers
these $2M-1$ coordinates from $M$ positive windows in one orientation
and $M-1$ in the other, all at two flights. Its inverse uses a
Vandermonde matrix with a shared constant coefficient; it does not
assume a normalization estimated at zero. The statistical orders are
regular parametric consequences of this physical inverse. The original
fixed-area, fixed-leading-hierarchy theorem and the general smooth
profile theorem retain their distinct hypotheses and conclusions.
